\exercisesfor{6}

\begin{enumerate}
\def\labelenumi{\arabic{enumi}.}
\item
  令 D 为 \(\mathbf{SL}(2,\mathbf{C})\) 中形如 \(\begin{pmatrix}a & b \\ 0 & a^{-1}\end{pmatrix}\) 的元素集合, 其中 a \textgreater{} 0 且 \(b \in C\). 证明 D 是 \(\mathbf{SL}(2,\mathbf{C})\) 的底层实李群的李子群, 并且映射 \((u,d) \mapsto ud\) 从 \(\mathbf{SU}(2,\mathbf{C}) \times \mathbf{D}\) 到 \(\mathbf{SL}(2,\mathbf{C})\) 是实解析流形的同构. 由此以及 §3 练习 7 (c) 推出 \(\mathbf{SL}(2,\mathbf{C})\) 是单连通的.
\item
  令 G 为 \(G\)\(\mathbf{SL}(2, \mathbf{R})\) 的万有覆盖, \(\pi\) 为从 G 到 \(G\)\(\mathbf{SL}(2, \mathbf{R})\) 的典范态射, N 为其核.\(N\)
\end{enumerate}

\begin{enumerate}
\def\labelenumi{(\alph{enumi})}
\item
  仿照练习 1, 证明存在从实解析流形 \(\mathbf{U} \times \mathbf{R}^2\) 到实解析流形 \(\mathbf{SL}(2, \mathbf{R})\) 的同构. 推出 N 同构于 \(N\)\(\mathbf{Z}\).
\item
  若 \(\rho\) 是 G 在复向量空间上的解析线性表示, 则 \(G\)\(\operatorname{Ker} \rho \supset N\). ( 的表示 \(\mathrm{L}(\rho)\) 通过复化定义出 \(\mathfrak{sl}(2, \mathbf{R})\)\(\mathfrak{sl}(2, \mathbf{C})\) 的一个表示, 后者根据练习 1 具有 \(\mathrm{L}(\sigma)\) 的形式, 其中 \(\sigma\) 是 \(\mathbf{SL}(2, \mathbf{C})\) 的一个解析线性表示. 然后 \(\mathrm{L}(\sigma \circ \pi) = \mathrm{L}(\rho)\), 因而 \(\sigma \circ \pi = \rho\).)
\end{enumerate}

\begin{enumerate}
\def\labelenumi{\arabic{enumi}.}
\setcounter{enumi}{2}
\item
  令 \(G_{1}\) 为 § 4 练习 5 中定义的单连通幂零实李群. 令 Z 为 \(G_{1}\) 的中心, N 为 Z 的非平凡离散子群, 且 \(G = G_{1}/N\). 若 \(\rho\) 是 G 的有限维解析线性表示, 则 \(\rho(Z/N)\) 是一组半单自同构, 因为 Z/N 是紧的; 但 \(Z = (G, G)\), 因而 \(\rho(Z/N)\) 是幂零的 (第 I 章 § 6 命题 6). 所以 \(\rho\) 在 Z/N 上是平凡的.
\item
  设 G 是紧连通复李群. 证明 G 的每个解析线性表示都是平凡的.
\item
  在 § 1 练习 9 的记号下, 证明 \(\mathbf{H}'\) 是 H 的积分子群. 推出在单连通群 \(\mathbf{SU}(2,\mathbf{C})\times \mathbf{SU}(2,\mathbf{C})\) (§ 3, 练习 7 (c)) 中存在不是闭的单参数子群.
\end{enumerate}

¶ 6. 令 I 为带离散拓扑的区间 \(\{1,2\}\). 令 E 为函数 \(f: I \to R\) 所成的完备赋范空间, 其中 \(\|f\| = \sum_{i \in I} |f(i)| < +\infty\). 对所有 \(x \in I\), 令 \(\varepsilon_x\) 为 E 中满足当  时 \(\varepsilon_x(t) = 0\) 且 \(t \neq x\)\(\varepsilon_x(x) = 1\) 的元素. 令 P 为由 \(x\varepsilon_x\) (\(x \in I\)) 生成的 E 的子群; 它是 E 的离散子群. 令 G 为李群 E/P. 令 F 为 E 中由满足  的 \(f \in E\) 组成的超平面. 令 H 为李群 F/(F \(\sum_{i \in I} f(i) = 0\) P). 令 \(\phi\) 为从 H 到 G 的典范态射. 证明 \(\phi\) 是双射且是浸入, 但不是李群同构. (令 f 为 E 上核为 F 的线性形式; 证明 \(f(P) = R\), 从而 \(F + P = E\).) 推出命题 3 中的可数性假设不可省略.

\begin{enumerate}
\def\labelenumi{\arabic{enumi}.}
\setcounter{enumi}{6}
\tightlist
\item
  令 \(\phi\) 为从 \(\mathbf{Z} / 2\mathbf{Z}\) 到 \(\mathbf{C}\) 的置换群的态射, 将非单位元映到 \(z\mapsto \overline{z}\). 令 \(\mathrm{G}\) 为由  相应的 \(\mathbf{Z} / 2\mathbf{Z}\) 与实李群 \(\mathbf{C}\) 构成的半直积. 它是一个李代数为 \(\phi\)\(\mathbf{R}^2\) 的实李群. 证明 \(\mathrm{G}\) 上不存在与其实李群结构相容的复李群结构.
\end{enumerate}

¶ 8. 设 H 是维数 \(\aleph_{0}\) 的复 Hilbert 空间, G 是视为实李群的 H 的酉群. 群 G 是单连通的. \(\dagger\) (a) 令 Z 为 G 的中心, 它同构于 T. 令 a 为无理数. 令 \(\hat{\mathbf{a}}\) 为 \(\mathrm{L}(\mathrm{G})\times \mathrm{L}(\mathrm{G})\) 中由 \((x,ax)\) (\(x\in \mathrm{L}(Z)\)) 组成的李子代数. 令 S 为相应的 \(\mathrm{G}\times \mathrm{G}\) 的积分子群. 则 \(\hat{\mathbf{a}}\) 是 \(\mathrm{L}(\mathrm{G})\times \mathrm{L}(\mathrm{G})\) 的理想; 然而 S 不是闭的, 且在 \(\mathbf{Z}\times \mathbf{Z}\) 中稠密.

\begin{enumerate}
\def\labelenumi{(\alph{enumi})}
\setcounter{enumi}{1}
\tightlist
\item
  令 \(\mathfrak{g} = (\mathrm{L}(\mathrm{G})\times \mathrm{L}(\mathrm{G})) / \tilde{\mathfrak{s}}\). 不存在李代数为 \(\mathfrak{g}\) 的李群. (若 \(\mathrm{H}\) 是这样的群. 由于 \(\mathrm{G}\) 是单连通的, 存在态射 \(\phi : \mathrm{G}\times \mathrm{G}\to \mathrm{H}\), 使 \(\mathrm{L}(\phi)\) 是从 \(\mathrm{L}(\mathrm{G})\times \mathrm{L}(\mathrm{G})\) 到 \(\mathfrak{g}\) 的典范映射. 令 \(\mathrm{N} = \operatorname{Ker}\phi\). 则 \(\mathrm{L}(\mathrm{N})\supset \tilde{\mathfrak{s}}\), 因而 \(\mathrm{N}\supset \mathrm{S}\), 再由 (a) 得 \(\mathrm{N}\supset \mathrm{Z}\times \mathrm{Z}\). 于是 \(\mathrm{L}(\mathrm{N})\supset \mathrm{L}(\mathrm{Z})\times \mathrm{L}(\mathrm{Z})\), 矛盾.)
\end{enumerate}

\begin{enumerate}
\def\labelenumi{\arabic{enumi}.}
\setcounter{enumi}{8}
\item
  设 X 是维数为 n 的非空连通紧复流形.~假设 X 上存在全纯向量场 \(\xi_{1},\ldots,\xi_{n}\), 它们在 X 的每一点都线性无关. 证明存在复李群 G 及其离散子群 D, 使 X 微分同胚于 G/D. (令 \(c_{ijk}\) 为 X 上满足 \([\xi_{i},\xi_{j}]=\sum_{k}c_{ijk}\xi_{k}\) 的全纯函数. 因为 X 紧, \(c_{ijk}\) 是常数. 取 G 为李代数以 \(c_{ijk}\) 为结构常数的单连通复李群, 并使用定理 5.)
\item
  设 G 是具有有限个连通分支的李群. G 是幺模的, 当且仅当对所有  都有 Tr ad \(a = 0\).\(a \in \mathrm{L}(\mathrm{G})\)
\item
  设 E 是 \(\mathbf{C}\) 上的完备可赋范空间, \(v \in \mathcal{L}(\mathrm{E})\), 且 \(g = \exp(v)\). 假设 \(\operatorname{Sp}(v) \cap 2i\pi(\mathbf{Z} - \{0\}) = \emptyset\). 令 \(\mathrm{E}_1, \mathrm{E}_2\) 为 E 的在 v 下不变的闭向量子空间, 且 \(v\)\(\mathrm{E}_2 \subset \mathrm{E}_1\). 假设由 g 在 \(\mathrm{E}_1 / \mathrm{E}_2\) 上导出的自同构是恒等映射. 则 \(g\)\(v(\mathrm{E}_1) \subset \mathrm{E}_2\).
\end{enumerate}

¶ 12. 考虑实或复完备可赋范空间 E, 以及 E 的一个闭向量子空间 F, 满足 \(F^{0}\) 在 \(E^{\prime}\dagger\) 中没有拓扑补空间 (a) 令 A (分别为 B) 为 E (分别为 F) 的连续自同态所成的完备可赋范空间. 令 C 为满足  的 \(u \in A\) 的集合. 令 \(u(F) \subset F\)\(\alpha\) 为从 C 到  的映射 \(u \mapsto (u, u \mid F)\). 则 \(A \times B\)\(\alpha\) 是从完备可赋范空间 C 到 \(A \times B\) 的一个闭向量子空间上的同构, 且 \(\alpha(C)\) 在 \(A \times B\) 中没有拓扑补空间. (令 \(x \in E\) 满足 \(x \notin F\). 令 \(\xi \in \mathrm{F}^0\) 满足 \(\langle x, \xi \rangle = 1\). 若 \(\eta \in \mathrm{E}'\), 令 \(\tilde{\eta}\) 表示 A 中的元素 \(y \mapsto \langle y, \eta \rangle x\). 假设存在从 A × B 到  的投影算子 \(\pi\). 定义 \(\alpha(\mathrm{C})\)\(\tilde{\pi} \colon \mathrm{E}' \to \mathrm{E}'\) 为 \(\tilde{\pi}(\eta) = t(\alpha^{-1}(\pi(\tilde{\eta}, 0)))(\xi)\). 则 \(\tilde{\pi}\) 是从 \(\mathrm{E}'\) 到 \(\mathrm{F}^0\) 的投影算子, 矛盾.)

如下: \(M = \bigoplus_{i=0}^{n} M_{i}\) 按 \(M_{i}\) 分次; \(M_{0}\) 是 1 的标量倍数的集合; \(M_{i}\) 是非零元素 u 的标量倍数的集合; \(M_{2} = \{0\}\); \(M_{3} = F\); \(M_{4} = E\); 当 i \(M_{i} = \{0\}\)\textgreater{} 5 时 ; 对所有 \(x \in M_{3}\) 有 ux = x. 令 N 为完备可赋范空间 M 的连续自同态所成的完备可赋范空间. 令 \(N_{1}\) 为 M 的连续导子的集合. 利用 (a) 证明 \(N_{1}\) 在 N 中没有拓扑补空间.

\begin{enumerate}
\def\labelenumi{(\alph{enumi})}
\setcounter{enumi}{2}
\tightlist
\item
  证明 M 的双连续自同构群是 \(\mathbf{GL}(\mathbf{M})\) 的李拟子群, 但不是 \(\mathbf{GL}(\mathbf{M})\) 的李子群. (使用命题 18 的推论 2 以及 (b).)
\end{enumerate}

\begin{enumerate}
\def\labelenumi{\arabic{enumi}.}
\setcounter{enumi}{12}
\item
  设 \(\mathbf{G}\) 是实李群, \(\mathbf{L} = \mathbf{L}(\mathbf{G})\), 且 \(\phi: \mathbf{L} \to \mathbf{G}\) 是在 0 处可微的映射, 满足 \(\mathrm{T}_0(\phi) = \mathrm{Id}_{\mathrm{L}}\), 并且对所有 \(\phi(nx) = \phi(x)^n\)\(x \in \mathbf{L}\) 和 \(n \in \mathbf{Z}\) 有 . 则 \(\phi = \exp_0\). (令 \(\mathbf{V}\) 为 \(\mathbf{L}\) 中 0 的一个邻域, \(\mathbf{W}\) 为 \(e\)\(\mathbf{G}\) 中 e 的一个邻域, 使得 \(0 = \exp_0 | \mathbf{V}\) 是从 \(\mathbf{V}\) 到 \(\mathbf{W}\) 的解析同构. 令 \(\psi = \theta^{-1} \circ (\phi | \phi^{-1}(\mathbf{W}))\). 则 \(\mathrm{T}_0(\psi) = \mathrm{Id}_{\mathrm{L}}\), 从而利用当 x 充分接近 0 且  时的等式 , 得 \(\psi = \mathrm{Id}_{\mathrm{L}}\).\(\psi\left(\frac{1}{n}x\right) = \frac{1}{n}\psi(x)\)\(x\)\(n \in \mathbf{Z} - \{0\}\)
\item
  令 \(\mathfrak{a}_1\) 为交换李代数 \(\mathbf{R}^4\), 将其识别为 \(\mathbf{C}^2\). 令 \(\mathfrak{a}_2\) 为交换李代数 \(\mathbf{R}\). 令 \(\phi\) 为从 \(\mathfrak{a}_2\) 到 \(\mathrm{Der}(\mathfrak{a}_1)\) 的同态, 将 1 映到导子 \((z_1, z_2) \mapsto (iz_1, i\sqrt{2}z_2)\). 令 \(\mathfrak{a}\) 为由  相应的 \(\mathfrak{a}_2\) 与 \(\mathfrak{a}_1\) 的半直积.\(\phi\)
\end{enumerate}

\begin{enumerate}
\def\labelenumi{(\alph{enumi})}
\tightlist
\item
  证明 (Int a) \textbar{} \(a_{1}\) 对所有 \(\phi \in R\) 都包含自同构
\end{enumerate}

\[
\left(z _ {1}, z _ {2}\right) \mapsto \left(e ^ {i \phi} z _ {1}, e ^ {i \sqrt {2} \phi} z _ {2}\right).
\]

\begin{enumerate}
\def\labelenumi{(\alph{enumi})}
\setcounter{enumi}{1}
\item
  证明 (Int a) \textbar{} a₁ 在 GL(a₁) 中的闭包 G 对所有 (φ, φ') ∈ R² 都包含自同构 (z₁, z₂) ↦ (eⁱΦz₁, eⁱΦ'z₂).
\item
  证明 \(\mathbf{L}(\mathbf{G})\) 包含  的自同态 \(u: (z_1, z_2) \mapsto (z_1, 0)\), 且不存在 \(\mathfrak{a}_1\)\(x \in \mathfrak{a}\) 使 \(u = (\operatorname{ad} x) \mid \mathfrak{a}_1\).
\item
  推出 \(\operatorname{Int}(\mathfrak{a})\) 在 \(\operatorname{Aut}(\mathfrak{a})\) 中不是闭的.
\end{enumerate}

\begin{enumerate}
\def\labelenumi{\arabic{enumi}.}
\setcounter{enumi}{14}
\tightlist
\item
  证明 §3 练习 7 (a) 的有限维单连通实李群包含非单连通的李子群 (事实上, 它包含同构于 U 的李群).
\end{enumerate}

  ¶ 16. (a) 设 g 是复李代数. 令 \(\bar{g}\) 为使用 C 的自同构 \(\lambda \mapsto \bar{\lambda}\) 从 g 导出的复李代数. 令 \(g_0\) 为实李代数
由限制标量从 g 导出. 令 \(g' \supset g_0\) 为复李代数 \(g_0 \otimes_{R} C\). 对所有 \(x \in g\), 记

\[
f (x) = \frac {1}{2} (x \otimes 1 - (i x) \otimes i) \in \mathfrak {g} ^ {\prime}, \quad g (x) = \frac {1}{2} (x \otimes 1 + (i x) \otimes i) \in \mathfrak {g} ^ {\prime}.
\]

证明 \(f(\text{resp. } g)\) 是从 \(\mathfrak{g}\) (分别从 \(\overline{\mathfrak{g}}\)) 到  的理想 \(\mathfrak{m}\) (分别为 \(\mathfrak{n}\)) 的同构, 且理想 \(\mathfrak{g}'\)\(\mathfrak{m}, \mathfrak{n}\) 在 \(\mathfrak{g}'\) 中互补. 这定义了从 \(\mathfrak{g}'\) 到 \(\mathfrak{m}, \mathfrak{n}\) 的投影算子. 证明对所有 \(x \in \mathfrak{g}, f(x) = p(x), g(x) = q(x)\), .

\begin{enumerate}
\def\labelenumi{(\alph{enumi})}
\setcounter{enumi}{1}
\item
  假设 g 是有限维的. 令 G 为李代数为 g 的单连通复李群. 令 \(\overline{G}\) 为共轭复李群, \(G_{0}\) 为底层实李群. 则 \(\mathrm{L}(\overline{\mathbf{G}})=\overline{\mathfrak{g}}\) 且 \(\mathrm{L}(G_{0})=g_{0}\). 令 M (分别为 N) 为李代数为 m (分别为 n) 的单连通复李群. 于是 f 定义了从 G 到 M 的同构 \(\phi\), g 定义了从  到 N 的同构 \(\gamma\), 而李代数为 \(\overline{G}\) 的单连通复李群 \(G'\) 与 M × N 对应. 证明 \(g'\)\(G'\) 中李代数为 \(g_{0}\) 的实积分子群是闭的、单连通的, 并与 \(G_{0}\) 对应. 将 \(pr_{1}\) (分别为 \(pr_{2}\)) 限制在 G  M × N 上, 得到从 G (分别从 ) 到 M (分别到 N) 的同构 \(\phi\) (分别为 \(\gamma\)).\(\overline{G}\)
\item
  由 (b) 以及第 10 号的注记 2 推出,  连同  到 \(\mathbf{G}'\) 的典范嵌入, 是 \(\mathbf{G}_0\) 的万有复化.\(\mathbf{G}'\)\(\mathbf{G}_0\)
\item
  设 H 是李代数为 g 的连通复李群, \(\bar{H}\) 是共轭复李群, \(H_{0}\) 是底层实李群, 且 G 是 H 的万有覆盖空间. 令 K 为 G 到 H 的典范态射的 (离散) 核. 证明 K 是 \(G'\) 的中心子群. 因而 \(G_{0}\) 到 \(G'\) 的典范嵌入定义了 \(H_{0}\) 到 \(G'/N\) 的嵌入 i. 证明 \((G'/N, i)\) 是 \(H_{0}\) 的万有复化. (注意这个有序对具有命题 20 的万有性质.)
\item
  记 \(G^{\prime}/N = (H_{0})_{\mathbf{c}}\). 由 (d) 推出从 \(\psi\)\((\mathrm{H}_{0})_{\mathbf{c}}\) 到 \(H \times \bar{H}\) 的典范态射 , 它是一个覆盖映射. 通过例子 \(H = C^{*}\) 证明 \(\psi\) 一般不是同构.
\end{enumerate}

\begin{enumerate}
\def\labelenumi{\arabic{enumi}.}
\setcounter{enumi}{16}
\tightlist
\item
  \begin{enumerate}
  \def\labelenumii{(\alph{enumii})}
  \tightlist
  \item
    令 G, \(\overline{G}\), \(G_{0}\) 如练习 16 (b) 中. 令 V 为有限维复向量空间, \(\rho\) 为 \(G_{0}\) 在 V 上的不可约线性表示. 证明存在有限维复向量空间 X, Y, G (分别为 ) 在 X (分别在 Y) 上的解析不可约线性表示 \(\sigma\) (分别为 \(\tau\)), 以及从 V 到 X \(\overline{G}\)\(\otimes\) Y 的同构, 将 \(\rho\) 变为 \(\sigma \otimes \tau\). (使用练习 16 (a)、第 I 章 § 2 的命题 2 以及《代数》第 VIII 章 § 7 的命题 8 的推论.)
  \end{enumerate}
\end{enumerate}

\begin{enumerate}
\def\labelenumi{(\alph{enumi})}
\setcounter{enumi}{1}
\tightlist
\item
  证明若去掉 G 单连通的假设, (a) 的结论不一定成立. (考虑复李群 \(\mathbf{C}^*\).)
\end{enumerate}

\begin{enumerate}
\def\labelenumi{\arabic{enumi}.}
\setcounter{enumi}{17}
\tightlist
\item
  设 \(\Lambda\) 是李代数为 \(a\) 的有限维连通交换复李群; 令 \(\Lambda\) 为 \(\exp_{A}\) 的核, 从而 \(\Lambda\) 与 \(a / \Lambda\) 对应.
\end{enumerate}

\begin{enumerate}
\def\labelenumi{(\alph{enumi})}
\tightlist
\item
  以下条件等价:
\end{enumerate}

(al) 典范映射 \(\mathbf{C} \otimes_{\mathbf{Z}} \Lambda \to a\) 是单射.

(a2) A 同构于某个 \((\mathbf{C}^{*})^{n}\) 的李子群.

(a3) A 同构于群 \((\mathbf{C}^{*})^{p}\times \mathbf{C}^{q}\).

(a4) A 具有有限维忠实复线性表示.

(a5) A 具有有限维半单忠实复线性表示, 且其像是闭的.

\begin{enumerate}
\def\labelenumi{(\alph{enumi})}
\setcounter{enumi}{1}
\tightlist
\item
  以下条件等价:
\end{enumerate}

(b1) 典范映射 \(\mathbf{C} \otimes_{\mathbf{Z}} \Lambda \to \mathfrak{a}\) 是满射.

(b2) A 同构于群 \((\mathbf{C}^{*})^{n}\) 的一个商.

(b3) A 的任何直因子都不同构于 C.

(b4) A 的每个复线性表示都是半单的.

\begin{enumerate}
\def\labelenumi{(\alph{enumi})}
\setcounter{enumi}{2}
\tightlist
\item
  以下条件等价:
\end{enumerate}

\begin{enumerate}
\def\labelenumi{(\roman{enumi})}
\setcounter{enumi}{149}
\tightlist
\item
  典范映射 \(\mathbf{C} \otimes_{\mathbf{Z}} \Lambda \to \mathfrak{a}\) 是双射.
\end{enumerate}

(c2) A 同构于群 \((\mathbf{C}^{*})^{n}\).

\begin{enumerate}
\def\labelenumi{(\alph{enumi})}
\setcounter{enumi}{3}
\tightlist
\item
  令 \(\mathbf{F}\) 为 \(\mathbf{A}\) 的有限子群, 令 \(\mathbf{A}' = \mathbf{A} / \mathbf{F}\). 证明 \(\mathbf{A}\) 满足条件 \((a_i)\) (分别为 \((b_i), (c_i)\)) 当且仅当 \(\mathbf{A}'\) 满足相应条件.
\end{enumerate}

\begin{enumerate}
\def\labelenumi{\arabic{enumi}.}
\setcounter{enumi}{18}
\tightlist
\item
  设 G 是实李群. 证明存在 L(G) 中 0 的一个邻域 V, 使得对 V 中的 \(G\)\(V\)\(L(G)\)\(x, y\) 有:\(V\)
\end{enumerate}

\[
\exp (t (x + y)) = \lim _ {n \rightarrow + \infty} \left(\exp \frac {t x}{n}. \exp \frac {t y}{n}\right) ^ {n}
\]

\[
\exp (t ^ {2} [ x, y ]) = \lim _ {n \rightarrow + \infty} \left(\exp \frac {t x}{n}. \exp \frac {t y}{n}. \exp \frac {- t x}{n}. \exp \frac {- t y}{n}\right) ^ {n ^ {2}}
\]

对 \(t \in (0, 1)\) 一致成立.

\begin{enumerate}
\def\labelenumi{\arabic{enumi}.}
\setcounter{enumi}{19}
\item
  设 G 是李群, H 是 G 的李子群, A 是 G 的积分子群, 且 \(\mathrm{L}(\mathrm{H}) \cap \mathrm{L}(\mathrm{A}) = \{0\}\). 证明 \(H \cap A\) 在李群 A 中是离散的.
\item
  设 \(G\) 是实李群, H 是正规的一维积分子群.\(H\)
\end{enumerate}

\begin{enumerate}
\def\labelenumi{(\alph{enumi})}
\item
  若 H 不是闭的, 则 \(\bar{H}\) 是紧的 (《谱理论》第 II 章 §2 引理 1), 因而同构于 \(T^{n}\); 若 G 连通, 则它在 G 中是中心的 (使用《一般拓扑》第 VII 章 §2 命题 5).
\item
  假设 H 是闭的. 令 a 为 G 的一个元素, C(a) 为 G 中与 a 交换的元素的集合. 假设 H ⊕ C(a).
\end{enumerate}

若 H 同构于 R, 则 \(\mathrm{C}(a) \cap \mathrm{H} = \{e\}\). (考虑 H 的自同构 \(\alpha: h \mapsto a^{-1}ha\).) 若 H 同构于 T, 则 \(\mathrm{C}(a) \cap \mathrm{H}\) 有两个元素. (再次考虑 \(\alpha\), 并使用《一般拓扑》第 VII 章 §2 命题 6.) 若 G 连通, 则第二种情形不可能发生 (使用 \((a)\)).

\begin{enumerate}
\def\labelenumi{(\alph{enumi})}
\setcounter{enumi}{2}
\tightlist
\item
  在 (b) 的假设下, 再假设 G/H 是交换的.
\end{enumerate}

于是 \(\mathrm{G} = \mathrm{C}(a)\) . H. (注意从 H 到 H 的映射 \(h \mapsto h^{-1}\alpha(h)\) 是满射.)

\begin{enumerate}
\def\labelenumi{\arabic{enumi}.}
\setcounter{enumi}{21}
\tightlist
\item
  设 G 是实李群, H 是正规的一维积分子群, A 是 G 的闭子群, 且 AH 在 G 中不是闭的. 则 H 是 AH 的单位元连通分支中的中心子群. (将问题归约为 G = AH 且 G 连通的情形. 令 B 为 A 中与 H 交换的元素所成的集合, 则 B 是 G 的正规子群. 对 B 取商, 将问题归约为 B = \(\{e\}\) 的情形. 由于 G 中每个换位子都与 H 交换, A 于是是交换的. 假设 H 是闭的并使用练习 21 (c), 证明 AH 将是闭的, 矛盾. 因而 H 不是闭的, 可以使用练习 21 (a).)
\end{enumerate}

¶ 23. 设 G 是实李群, \(c = (\mathrm{U}, \phi, \mathrm{E})\) 是以单位元 e 为中心的 G 的流形上的一个图. 若 \(e\)\(x \in \mathrm{U}\), 令 \(|x|_c\) 表示 Banach 空间 E 中 \(\phi(x)\) 的范数.

\begin{enumerate}
\def\labelenumi{(\alph{enumi})}
\tightlist
\item
  证明若 \(c' = (\mathbf{U}', \phi', \mathrm{E}')\) 是以 e 为中心的另一个图, 则存在常数 \(e\)\(\lambda > 0\) 和 \(\mu > 0\), 使得
\end{enumerate}

\[
\mu | x | _ {c} \leqslant | x | _ {c ^ {\prime}} \leqslant \lambda | x | _ {c}
\]

对所有充分接近 e 的 \(x \in G\) 都成立.

\begin{enumerate}
\def\labelenumi{(\alph{enumi})}
\setcounter{enumi}{1}
\tightlist
\item
  证明对所有 \(\rho > 0\), 存在包含于  的 e 的邻域 \(\mathbf{U}_{\rho}\), 使得对 \(e\) 中的 \(\mathbf{U}\)\(x, y\), \(\mathbf{U}_{\rho}\)\((x, y) \in \mathbf{U}_{\rho}\), 且
\end{enumerate}

\[
\left| (x, y) \right| _ {c} \leqslant \rho . \inf (\left| x \right| _ {c}, \left| y \right| _ {c}).
\]

\begin{enumerate}
\def\labelenumi{(\alph{enumi})}
\setcounter{enumi}{2}
\tightlist
\item
  假设 \(\mathbf{G}\) 是有限维的. 令 \(\Gamma\) 为 \(\mathbf{G}\) 的离散子群. 对 \(0 < \rho < 1\) 使用 (b), 并选取相对紧的 \(\mathrm{U}_{\rho}\). 集合 \(\mathrm{U}_{\rho} \cap \Gamma\) 是有限的. 令
\end{enumerate}

\[
\{e = \gamma_ {0}, \gamma_ {1}, \dots , \gamma_ {m} \}
\]

是它的元素, 并按如下方式编号:

\[
\left| \gamma_ {0} \right| _ {c} \leqslant \left| \gamma_ {1} \right| _ {c} \leqslant \dots \leqslant \left| \gamma_ {m} \right| _ {c}.
\]

证明若 \(i \leqslant m\) 且 \(j \leqslant m\), 则换位子 \((\gamma_{i}, \gamma_{j})\) 等于某个  的 \(\gamma_{k}\). 推出由 \(k < \inf(i,j)\)\(\mathrm{U}_{\rho} \cap \Gamma\) 生成的子群是类 \(\leqslant m\) 的幂零群.

由上述结果推出存在 e 的一个邻域 \(\mathbf{V}\), 使得对每个 \(e\) 的离散子群 \(\Gamma\), 都存在包含  的 \(\mathbf{G}\) 的幂零积分子群 \(\mathbf{N}\).\(\mathbf{G}\)\(\mathbf{V} \cap \Gamma\)

\begin{enumerate}
\def\labelenumi{(\alph{enumi})}
\setcounter{enumi}{3}
\tightlist
\item
  假设 G 是连通且有限维的, 并包含离散子群的递增序列 \(D_{n}\), 其并在 G 中稠密. 则 G 是幂零的. (利用 (c), 证明存在一个中心的单参数子群 H, 它与充分大的 n 所对应的 \(D_{n}\) 相交于一个不同于 e 的点. 按  归纳, 根据 H 是闭的还是相对紧的区分两种情形 (练习 21 (a)).
\end{enumerate}

\begin{enumerate}
\def\labelenumi{\arabic{enumi}.}
\setcounter{enumi}{23}
\tightlist
\item
  设 \(G, G'\) 是有限维实李群, f 是从 G 到 \(f\)\(G\)\(G'\) 的满射态射, N 是其核, H 是 G 的积分子群, \(N\)\(H\)\(G\)\(H' = f(H)\).
\end{enumerate}

\begin{enumerate}
\def\labelenumi{(\alph{enumi})}
\item
  假设 N 是有限的. \(H'\) 是闭的, 当且仅当 H 是闭的.
\item
  假设 N 是紧的. 若 H 是闭的, 则 \(H'\) 是闭的.
\end{enumerate}

\begin{enumerate}
\def\labelenumi{\arabic{enumi}.}
\setcounter{enumi}{24}
\tightlist
\item
  设 \(G\) 是有限维实或复李群. 令 \(S \subset L(G)\). 假设 L(G) 是由 S 生成的李代数, 且 S 在伸缩变换下不变.\(L(G)\)\(S\)\(S\)
\end{enumerate}

\begin{enumerate}
\def\labelenumi{(\alph{enumi})}
\item
  令 H 为由 exp S 生成的 G 的子群. 证明 H 在 G 中是开的. (令 A 为满足  的 \(x \in L(G)\) 的集合, B 为由 A 生成的 L(G) 的向量子空间. 证明 \(\exp(Kx) \subset H\)\(L(G)\){[}A, S{]} \(\subset A\), 从而 {[}A, A{]} \(\subset B\). 推出 B 是 L(G) 的子代数, 然后使用 § 4 的命题 3.)\(L(G)\)
\item
  再假设
\end{enumerate}

\[
(x \in S \text {   and   } y \in S) \Rightarrow ((\operatorname{Ad} \exp x) (y) \in S).
\]

则由  生成的  的向量子空间 \(\mathbf{V}\) 等于 \(\mathbf{L}(\mathbf{G})\)\(\mathbf{S}\)\(\mathbf{L}(\mathbf{G})\). (利用 \((a)\), 证明 \([\mathbf{L}(\mathbf{G}), \mathbf{V}] \subset \mathbf{V}\).)

\begin{enumerate}
\def\labelenumi{\arabic{enumi}.}
\setcounter{enumi}{25}
\tightlist
\item
  设 G 是 n 维连通紧复李群, X 是有限维连通复解析流形, \((g, x) \mapsto gx\) 是 G 在 X 上的解析左作用法则. 对所有 \(g \in G\), 令 \(\rho(g)\) 为从 X 到 X 的映射 \(x \mapsto gx\). 假设对所有不同于 e 的 \(g \in G\), \(e, \rho(g) \neq Id_x\). 则 G 在 X 中的每个轨道都是 X 的 n 维闭子流形. (令 \(x \in X\), H 为 x 的稳定子, H' 为 H 的单位元连通分支. 对所有 \(H'\)\(g \in H'\), 令 \(u(g)\) 为 \(\rho(g)\) 在 x 处的切映射. 仿照第 3 号命题 6 的论证, 证明对所有 , \(u(g)\) 都是恒等映射. 使用 § 1 练习 4 以及 X 的连通性, 推出 \(g \in H'\)\(H' = \{e\}\).)
\end{enumerate}

¶ 27. 设 G 是紧实李群, M 是  类紧流形, J 是包含 0 的 R 的一个开区间. 令 \((s, \langle x, \xi \rangle) \mapsto (m_{\xi}(s, x), \xi)\) 为从 G × (M × J) 到 M × J 的  类映射, G 通过该映射在 M × J 上左作用.

\begin{enumerate}
\def\labelenumi{(\alph{enumi})}
\item
  令 X 是 M × J 上 \(C^{1}\) 类向量场, 使得对所有 \(M \times J\)\((x, \xi) \in M \times J\), \(\mathbf{X}_{(x, \xi)}\) 的第二投影是 J 的切向量 l. 将 X 用 G 平移并在 G 上积分, 推出存在一个与 X 具有相同性质且进一步在 G 下不变的场 \(X'\).
\item
  证明存在从  到自身的微分同胚 \((x, \xi) \mapsto (h_{\xi}(x), \xi)\), 使得\(\mathbf{M} \times \mathbf{J}\)
\item
  对所有 \(\xi \in \mathbf{J}\), \(h_{\xi}\) 都是从 \(\mathbf{M}\) 到 \(\mathbf{M}\) 的微分同胚;
\end{enumerate}

\begin{enumerate}
\def\labelenumi{(\roman{enumi})}
\setcounter{enumi}{1}
\tightlist
\item
  \(m_{\xi}(s,x) = h_{\xi}(m_0(s,h_{\xi}^{-1}(x)))\) 对 \(s\in \mathbf{G}, x\in \mathbf{M},\xi \in \mathbf{J}\) 成立.
\end{enumerate}

(使用 (a) 和第 8 号定理 5.)

\begin{enumerate}
\def\labelenumi{\arabic{enumi}.}
\setcounter{enumi}{27}
\item
  举出一个有限维实李群 G 以及 G 的两个积分子群 A, B 的例子, 使得 \(\mathrm{L}(\mathrm{G})=\mathrm{L}(\mathrm{A})+\mathrm{L}(\mathrm{B})\), 但 AB 不是 G 的子群 (参见《积分论》第 VII 章 § 3 第 3 号命题 6).
\item
  设 G 是有限维连通复李群, \(G_{0}\) 是底层实李群, H 是 \(G_{0}\) 的积分子群. 证明存在包含 H 的 G 的最小积分子群 \(H^{*}\). 举出 H 在 \(G_{0}\) 中闭而 \(H^{*}\) 在 G 中不闭的例子 (取 \(G = C^{2}/Z^{2}\)).
\item
  设 \(G\) 是紧连通复李群, 因而形如 \(\mathbf{C}^n / D\), 其中 D 是 \(D\)\(\mathbf{C}^n\) 的秩为 \(2n\) 的离散子群.
\end{enumerate}

\begin{enumerate}
\def\labelenumi{(\alph{enumi})}
\item
  证明 G 上每个全纯 1 次微分形式 \(\omega\) 都是不变的. (令 \(\pi: \mathbf{C}^n \to \mathbf{G}\) 为典范态射. 令 \(\zeta_1, \ldots, \zeta_n\) 为 \(\mathbf{C}^n\) 上的坐标函数. 则 \(\pi^*(\omega)\) 形如 \(\sum_j a_j d\zeta_{j}\), 其中 \(a_j\) 是 \(\mathbf{C}^n\) 上在 D 下不变的全纯函数, 因而是常数.)
\item
  令 \(G' = C^{n}/D'\), 其中 \(D'\) 是 \(C^{n}\) 的秩为 2n 的离散子群. 证明每个解析流形同构 \(u: G \to G'\) 都具有形式 \(s \mapsto v(s) + a'\), 其中 v 是李群同构且 \(a' \in G'\). (令 \(\pi': C^{n} \to G'\) 为典范态射. 存在解析流形同构 \(\tilde{u}: C^{n} \to C^{n}\), 使 \(u \circ \pi' \circ \tilde{u}\). 对 G' 上的每个全纯 1 次微分形式 \(\omega'\), \(G'\)\(u^*(\pi^*(\omega'))\) 是 \(C^{n}\) 上的 1-形式, 根据 (a) 在平移下不变. 推出 \(\tilde{u}\) 是仿射映射.)
\end{enumerate}

